% This is samplepaper.tex, a sample chapter demonstrating the
% LLNCS macro package for Springer Computer Science proceedings;
% Version 2.21 of 2022/01/12
%
\documentclass[runningheads,10pt]{llncs}
%
\usepackage[T1]{fontenc}
% T1 fonts will be used to generate the final print and online PDFs,
% so please use T1 fonts in your manuscript whenever possible.
% Other font encondings may result in incorrect characters.
%
\usepackage{graphicx}
\usepackage{hyperref}
% Used for displaying a sample figure. If possible, figure files should
% be included in EPS format.
%
% If you use the hyperref package, please uncomment the following two lines
% to display URLs in blue roman font according to Springer's eBook style:
%\usepackage{color}
%\renewcommand\UrlFont{\color{blue}\rmfamily}
%\urlstyle{rm}
%
\begin{document}
%
\title{PT-AML: Machine learning framework to identified personalized treatments for acute myeloid leukemia}

\titlerunning{PT-AML}
% If the paper title is too long for the running head, you can set
% an abbreviated paper title here
%
\author{Siddhi P. Jani\inst{1} \and
Raghvendra Mall\inst{2,3}}
%
\authorrunning{Jani et al.}
% First names are abbreviated in the running head.
% If there are more than two authors, 'et al.' is used.
%
\institute{Centre of Brain Research, Indian Institute of Sciences, Bangalore, Karnataka, India
\email{siddhijani@cbr-iisc.ac.in}\\
\and
Biotechnology Research Center, Technology Innovation Institute, P.O. Box 9639, Abu Dhabi, UAE  \and
Department of Immunology, St. Jude Children's Research Hospital, Memphis, TN 38105, USA\\
\email{raghvendra.mall@tii.ae}}
%
\maketitle              % typeset the header of the contribution
%
\begin{abstract}
Acute myeloid leukemia (AML) is a cancer of myeloid cells with limited treatment options. 
The BeatAML cohort collected and characterized samples of $805$ patients over 10 years with the integration of \textit{ex vivo} drug sensitivity, clinical annotations, and DNA and RNA sequencing. 
Through association analysis, strong cross-cohort concordance and features of drug response were previously identified. 
In this work, we propose PT-AML, which leverages the BeatAML cohort and their collection strategy to devise and validate machine learning techniques for predicting drug response of AML patients given their clinical, genomic and transcriptomic profiles.
Our framework engineers vector representations for drugs, their relevant pathways, patient profiles and their cell states to accurately estimate drug response. 
The optimal PT-AML model attains a cross-validation mean-absolute-error (MAE) and spearman correlation (r) of $21.02 \pm 1.5$, $0.85 \pm 0.02$ respectively and MAE of 39.04 and r of 0.66 on the unseen test set respectively highlighting the generalization capability of the model.  

\keywords{Machine Learning \and AML \and Precision Medicine \and Cell States \and Drug Response}
\end{abstract}

\section{Introduction}
Acute Myeloid Leukemia (AML) is a highly heterogeneous malignancy characterized by diverse genetic and molecular features \cite{cancer2013genomic,papaemmanuil2016genomic}, presenting a significant challenge in effective treatment \cite{pulte2016survival}. Despite advancements in targeted therapies, predicting drug response for individual patients remains a major obstacle \cite{tyner2018functional}. Precision medicine initiatives such as the BeatAML project aim to address this challenge by integrating multi-omics data with clinical information to tailor therapies to individual patients \cite{tyner2018functional,bottomly2022integrative,burd2020precision}.

The BeatAML cohort, led by the Leukemia \& Lymphoma Society (LLS), is a precision medicine initiative designed to improve treatment for AML patients by matching them to targeted therapies based on their genetic profiles \cite{burd2020precision}. 
This study presents a harmonized dataset of 942 specimens from 805 AML patients collected in tranches. 
%The first tranche of patient sample data (denoted as waves 1 + 2) consists of 649 specimens from 526 patients and forms our training set \cite{tyner2018functional}. 
%An additional tranche representing the final two waves (waves 3 + 4) comprise a total of 293 specimens from 279 patients and corresponds to the test set \cite{bottomly2022integrative}. 
In \cite{bottomly2022integrative}, the authors demonstrated strong cross-cohort concordance in mutations and gene expression through association analysis. 
Through \textit{ex vivo} drug sensitivity testing, genomic and transcriptomic analysis, and clinical annotations, the study \cite{bottomly2022integrative} explores the relationship between drug response and AML cell differentiation states, uncovering significant correlations between gene expression and drug sensitivity. 
The study \cite{bottomly2022integrative} further emphasized the importance of understanding how drug response is influenced by both genetic mutations and the differentiation state of AML cells. 
It highlighted the interactions between cell type and other patient features can significantly affect the drug response. Despite these significant findings, there are several limitations including the retrospective nature, potential selection bias, and owing to just association analysis the inability to predict drug response for new patients. 

Thus, our work focuses on leveraging the rich multi-omics and clinical data from the BeatAML cohort to develop machine learning models \cite{lecun2015deep} which will provide personalized treatments for AML patients (\textbf{PT-AML}). PT-AML model will accurately predict drug responses based on patients' unique molecular and clinical profiles.  
Specifically, we build predictive models that integrate transcriptomic, genomic, and clinical features to understand the intricate relationship between disease state and drug response. 
By developing this model, we aim to identify the most effective drug for a patient given their multi-omics information.
%%%%%%%%%%%%%%%%%%%%%%%%%%%%%%%%%%%%%%%%%%%%%%%%%%%%%
\section{Materials and Methods}
\subsection{Data Collection and Processing} 
We utilized data from the BeatAML cohort \cite{bottomly2022integrative}. The cohort is divided into two primary groups: Waves 1 + 2 and Waves 3 + 4. Waves 1 + 2 include 649 specimens from 526 patients, while Waves 3 + 4 comprise 293 specimens from 279 patients, resulting in a total of 942 samples from 805 patients. For our analysis, we focused on $560$ samples, with 372 samples from Waves 1 + 2 forming our training set and $188$ from Waves 3 + 4 forming our test set. This is because mutation, transcriptomics, drug response and clinical data were available consistently for all these samples. Clinical data for these samples were curated, including age, sex, disease status (alive/dead), and clinical traits such as relapse, de novo AML, transformed AML, and prior malignancies were also included. The age of patients ranged from 1 to 88 years, with 245 female and 315 male patients. %A total of 217 patients were alive, 321 had passed, and the status for 22 patients was unknown.
%%%%%%%%%%%%%%%%%%%%%%%%%%%%%%%%%%%
\subsection{RNA-Seq \& DNA-Seq Analysis}
We processed RNA-seq data from both Waves 1 + 2 and Waves 3 + 4. From the clinical data, samples with matching gene expression data were retained, and all others were excluded. We identified 150 highly variable genes and oncogenes, whose scaled expressions were subsequently used as gene expression profiles. Clinical annotations were merged with the gene expression data. We next identified genesets representing functional processes for individual patients. The genesets comprising cell states such as cDC-like, GMP-like, HSC-like, Monocyte-like, Progenitor-like, and Promono-like and genes which are part of modules were obtained from the original BeatAML study \cite{bottomly2022integrative}. Similarly, genesets forming oncogenic pathways were obtained from \cite{roelands2020oncogenic}. 

We then performed oncogenic pathway activity \cite{roelands2020oncogenic,mall2021network}, module and cell state \cite{bottomly2022integrative} enrichment, estimation by undertaking single sample gene set enrichment using the GSVA (v1.50.5) package \cite{hanzelmann2013gsva} in R (v4.4.1). These scores highlight how active or enriched are these functional processes in individual patients. We used the normalized scores ([-0.5,0.5]) for pathway activity and module enrichment where higher scores correspond to higher activity. In the case of cell states, we used the un-normalized enrichment score ([0,1] as each sample would show significant enrichment for a single cell-state. For the DNA-Seq analysis, somatic variant calls were also incorporated to evaluate mutation frequencies and types. The frequency of mutation in top mutated genes and their corresponding mutation types were used as features from DNA-Seq data. Figure \ref{fig:fig1} highlights these characteristics of patient samples as a Heatmap.
\begin{figure}[!ht]
    \centering
    \includegraphics[width=1.1\textwidth]{figures/Figure1.pdf}
    \caption{Comprehensive Multi-omics characterization of patient samples along with clinical annotations.}
    \label{fig:fig1}
\end{figure}
%%%%%%%%%%%%%%%%%%%%%%%%%%%%%%%%%%%
\subsection{Drug Data and Fingerprint Generation} 
We collected drug response data from the BeatAML dataset. For each drug, Morgan molecular fingerprints (MFP) \cite{yang2022concepts} were derived using RDKit, and SMILES-based encodings of drugs were generated using a Teacher Forcing - Long Short Term Memory (TF-LSTM) deep learning model from \cite{mall2021modeling}. In our machine learning models, we used either the MFP or the encoding representation as a feature vector for drugs. 
%%%%%%%%%%%%%%%%%%%%%%%%%%%%%%%%%%
\subsection{Drug Pathway Distance Estimation}
To estimate an approximate computational distance of drugs from individual oncogenic pathways, we utilize the drug target (gene) information available via DrugBank \cite{knox2024drugbank}. Here we use the gene and corresponding protein as one or interchangeably. We then utilize the interaction network of proteins which is obtained through the \href{https://string-db.org/}{STRING} protein-protein interaction (PPI) database. From the STRING PPI network, we consider the network of human proteins and only those interactions which have strong experimental evidence. 

Considering the drug targets as starting points in this interaction network, we perform network propagation using a random-walk with restart approach \cite{sherif2022immune,mall2013kernel} to obtain a propagation vector for each drug. The propagation vector can be used as a surrogate for the global distance of a protein from a given drug (and its targets). To obtain a local propagation vector i.e. patient specific propagation vector, we multiply the propagation score of each protein with their corresponding gene expression available via RNA-Seq for that patient. 

Our aim is to estimate the proximity of an oncogenic pathway from a given drug and a patient profile. If the genes which are part of an oncogenic pathway are all ranked higher in the sorted propagation vector then the drug has close proximity to that oncogenic pathway where as if the genes are ranked randomly in the sorted propagation vector then the drug has low proximity to the oncogenic pathway. This can be estimated using the AUCell package (v1.24.0) \cite{aibar2017scenic} in R. It estimates the area under the receiver operating curve (AUC), where higher AUC corresponds to closer proximity and lower AUC corresponds to further distance between a drug and the oncogenic pathway for a given patient. Figure \ref{fig:fig2} highlights the steps involved in drug pathway distance estimation. 
\begin{figure}[!ht]
    \centering
    \includegraphics[width=0.99\textwidth]{figures/Figure2.pdf}
    \caption{Steps involved in patient-specific drug-pathway distance estimation.}
    \label{fig:fig2}
\end{figure}
%effect of drugs on specific pathways, we used a random walk with restart approach on STRING PPI networks. Pathway activity was estimated using gene expression profiles combined with diffusion scores from drug targets. The AUCell package was employed to measure the proximity of each pathway to the drugs in question. Drugs that had the highest impact on pathways were ranked, and their influence on disease outcomes was assessed.
%%%%%%%%%%%%%%%%%%%%%%%%%%%%%%%%%%%%
\subsection{Machine Learning Models}
We developed and trained multiple machine learning models to predict drug response for patients using combined molecular, genomic, and clinical data together with the drug and functional process representations. 
The designed models included generalized linear regression (GLR) \cite{nelder1972generalized}, support vector regression (SVR) \cite{chandorkar2015fixed}, Random Forests (RF) \cite{breiman2001random}, XGBoost \cite{chen2016xgboost}, LightGBM \cite{ke2017lightgbm}, CatBoost \cite{prokhorenkova2018catboost}, feed forward neural networks \cite{hastie2009elements} among the standard machine learning models. We also trained several deep learning models using the SMILES representation of drugs as input to convolutional neural networks \cite{sermanet2012convolutional}, LSTM \cite{gers2000learning} and graph attention networks \cite{velivckovic2017graph} combined with patient information and feed forward layers to build deep learning based drug response models. 

All models were trained using 5-fold cross-validation on the training set (Waves 1 + 2) and tested on Waves 3 + 4. Performance metrics, including mean absolute error (MAE), root mean squared error (RMSE), pearson correlation (r), spearman correlation (sr) and coefficient of determination (R$^2$) were computed to assess model performance. Lower MAE and RMSE values are better and higher r, sr and R$^2$ values are better.  All the tree-based models achieved near optimal performance with the Catboost model achieving the lowest test MAE. Figure \ref{fig:fig3} represents all the steps involved in building the PT-AML model.
\begin{figure}[!ht]
    \centering
    \includegraphics[width=0.99\textwidth]{figures/Figure3.pdf}
    \caption{The workflow followed to build the proposed PT-AML model.}
    \label{fig:fig3}
\end{figure}
%%%%%%%%%%%%%%%%%%%%%%%%%%%%%%%%%%%%%%%
\section{Results}
\subsection{Benchmarking Machine Learning Models}
We highlight the performance of different machine learning models in Figure \ref{fig:fig4}.
Figure \ref{fig:fig4}A highlights the distribution of the area under the drug response curve (AUC) between patient profiles in Wave 1+2 (Y\_train) and Wave 3+4 (Y\_test). We observe that the AUC values range between 0 and 300 where higher values indicate better response. Moreover, from Figure \ref{fig:fig4}A, we observe that the distribution of the drug response are similar between the training and test (with no statistical difference i.e. p-value $> 0.05$ based on Wilcoxon-ranksum test \cite{mann1947test}).

\begin{figure}[!ht]
    \centering
    \includegraphics[width=0.95\textwidth]{figures/Figure4.pdf}
    \caption{Performance comparison of machine learning models for predicting drug response of AML patients. We highlight the top 10 drugs for which our optimal model performs best and top 10 patient profiles for which the prediction was most accurate. Here `MFP' represents morgan fingerprints, `LS' corresponds to TF-LSTM embeddings and `SMILES' is the SMILES representation for drugs. Additionally `Feat' corresponds to the multi-omic feature representation for patients.}
    \label{fig:fig4}
\end{figure}
We then showcase the performance of the different machine learning and deep learning models in Figures \ref{fig:fig4}B-K and Table \ref{table:t1}. Using different representations for drugs namely the Morgan fingerprint (MFP) and TF-LSTM based embeddings (LS), we have two model for each standard machine learning technique like GLR, SVR, RF, XGBoost, CatBoost and LightGBM respectively. From Figures \ref{fig:fig4}B,C,H and Table \ref{table:t1}, we observe that the linear models and the feed forward neural network (FFNN) tend to be worst w.r.t. MAE and correlation metrics (r, sr, R$^2$) when compared to other methods. While linear models cannot capture the inherent non-linearity in the data, the FFNN achieves the best cross-validation scores across all models (see Table \ref{table:t1}) but has poor generalization performance indicating that the model has overfitted on the training data.

From Figure \ref{fig:fig4}D-G, we observe that the tree-based models achieve an MAE of $\approx$ 39 and a pearson correlation of $\approx$ 0.66 for the test set. The embedding based representation (LS) for drugs is slightly better than the MFP representation of drugs for RF, XGBoost and LightGBM as observed in Table \ref{table:t1}. The RF model (LS + Feat) achieves a cross-validation r and sr of 0.89 $\pm$ 0.01, R$^2$ of 0.79 $\pm$ 0.02 as depicted in Table \ref{table:t1}. Similarly, the LigthGBM model (LS + Feat) attains a cross-validation r of 0.89 $\pm$ 0.02, sr of 0.87 $\pm$ 0.02, and R$^2$ of 0.79 $\pm$ 0.03 as highlighted in Table \ref{table:t1}. 

However, the Catboost model using MFP representation for drugs has the best MAE of 39.04, RMSE of 52.03, r of 0.66 and R$^2$ of 0.44 for the test set as depicted in Table \ref{table:t1}. Interestingly, we observe that deep learning based methods tend to perform better than linear methods but worse than tree-based models (see Table \ref{table:t1}). The GAT, CNN, LSTM model in combination with FFNN achieve MAE of 40.7, 40.85, 40.7 and r of 0.63, 0.63, 0.62 respectively on the test set as illustrated in Figures \ref{fig:fig4}I-K and Table \ref{table:t1}. 

Overall the tree-based models (Catboost in particular) can best capture the intricate relationship between drug representations, multi-omic patient profile and functional processes to accurately estimate the drug reponse for given patient as depicted by their cross-validation performance and the generalization performance on independent test set.
%%%%%%%%%%%%%%%%%%%%%%%%%%%%%%%%%%%%%%%%%%%%%%%
\subsection{Top Patient and Drug Predictions}
Figure \ref{fig:fig4}L highlights the top 10 best patient profiles for which the PT-AML (hereby Catboost) model has the best performance w.r.t. the pearson correlation (r) and R$^2$ scores. These scores are estimated for each patient based on the total number of drugs tested (AUC scores in test set) and can vary per patient as indicated in `brackets' in Figure \ref{fig:fig4}L. We observe from Figure \ref{fig:fig4}L (top panel) that for each of these patients the r and R$^2$ values were > 0.7, indicating a strong alignment between the predicted and actual drug response. Similarly, we highlight the bottom 10 worst patient profiles in Figure \ref{fig:fig4}M (top panel). The correlation between the predicted and actual AUC scores for the worst patient (BA3117R) was still attaining a correlation (r) $\approx$ 0.2. It is noteworthy that for each patient there is more than 100 drug response curves available in the test set as depicted in Figures \ref{fig:fig4}L and M. This can enable the machine learning models to better capture patient specific multi-omic information driving the drug response (sensitivity vs resistance).

Figure \ref{fig:fig4}L (bottom panel) indicates the top 10 best drug predictions for which the PT-AML model has the best performance w.r.t. r and R$^2$ scores. These scores are estimated for each drug based on the total number of patient samples tested (AUC scores in test set) and vary per drug as indicated in `brackets' in Figure \ref{fig:fig4}L. We observe from Figure \ref{fig:fig4}L (bottom panel) that for drugs such as Ventoclax, Selumetinib, and Trametinib, the Catboost model can achieve correlation (r) values close to 0.7 and above, demonstrating accurate predictive performance. For each of the top 10 best drug predictions, there were over 100 patient samples available in the test set (see Figure \ref{fig:fig4}L). However, from Figure \ref{fig:fig4}M (bottom panel), we observe that there are certain drugs such as Perhexiline maleate, Baiclein, AF101, Ranolazaine, etc. the model predictions are fairly inaccurate resulting in poor correlation scores. This can partly be explained owing to lack of enough patient profiles for which these drugs have response curves as majority of these drugs have been tested for < 100 patient samples. 

Overall these results emphasize that while PT-AML model is effective for certain patients and drugs, challenges persist while predicting response for others. 
\begin{table}[!ht]
    \centering
    \resizebox{0.99\textwidth}{!}{
    \begin{tabular}{|c|c|c|c|c|c|c|c|c|c|c|c|}
    \hline
    \textbf{Methods} & \textbf{Features} & \textbf{CV (MAE)} & \textbf{Test (MAE)} & \textbf{CV (RMSE)} & \textbf{Test (RMSE)} & \textbf{CV (R$^2$)} & \textbf{Test (R$^2$)} & \textbf{CV (r)} & \textbf{Test (r)} & \textbf{CV (sr)} & \textbf{Test (sr)} \\
    \hline
    GLR & LS + Feat & 30.22 $\pm$ 1.97 & 44.86 & 40.34 $\pm$ 2.91 & 60 & 0.56 $\pm$ 0.05 & 0.28 & 0.74 $\pm$ 0.03 & 0.53 & 0.74 $\pm$ 0.03 & 0.56 \\
    GLR & MFP + Feat & 29.93 $\pm$ 1.92 & 43.48 & 40.10 $\pm$ 2.89 & 57.26 & 0.56 $\pm$ 0.05 & 0.34 & 0.75 $\pm$ 0.03 & 0.58 & 0.75 $\pm$ 0.03 & 0.58 \\
    SVR & LS + Feat & 30.13 $\pm$ 1.96 & 44.97 & 40.25 $\pm$ 2.90 & 60.37 & 0.56 $\pm$ 0.05 & 0.28 & 0.75 $\pm$ 0.03 & 0.53 & 0.74 $\pm$ 0.03 & 0.56 \\
    SVR & MFP + Feat & 29.93 $\pm$ 1.92 & 43.89 & 40.09 $\pm$ 2.89 & 57.73 & 0.56 $\pm$ 0.05 & 0.33 & 0.75 $\pm$ 0.03 & 0.57 & 0.75 $\pm$ 0.03 & 0.57 \\
    RF & LS + Feat & 21.82 $\pm$ 0.97 & 39.63 & 29.86 $\pm$ 1.51 & 52.08 & 0.79 $\pm$ 0.02 & 0.44 & 0.89 $\pm$ 0.01 & 0.66 & 0.89 $\pm$ 0.01 & 0.66 \\
    RF & MFP + Feat & 20.75 $\pm$ 1.22 & 40.16 & 28.29 $\pm$ 1.83 & 52.53 & 0.82 $\pm$ 0.02 & 0.44 & 0.90 $\pm$ 0.01 & 0.66 & 0.9 $\pm$ 0.01 & 0.65 \\
    XGBoost & LS + Feat & 24.29 $\pm$ 1.21 & 39.84 & 31.94 $\pm$ 1.76 & 52.34 & 0.73 $\pm$ 0.04 & 0.43 & 0.85 $\pm$ 0.02 & 0.66 & 0.84 $\pm$ 0.02 & 0.65 \\
    XGBoost & MFP + Feat & 22.94 $\pm$ 1.01 & 40.53 & 30.10 $\pm$ 1.45 & 52.88 & 0.78 $\pm$ 0.04 & 0.43 & 0.88 $\pm$ 0.2 & 0.66 & 0.87 $\pm$ 0.02 & 0.64 \\
    LightGBM & LS + Feat & 21.53 $\pm$ 1.12 & 39.5 & 28.75 $\pm$ 1.63 & 52.09 & 0.79 $\pm$ 0.03 & 0.44 & 0.89 $\pm$ 0.02 & 0.66 & 0.87 $\pm$ 0.02 & 0.66 \\
    LightGBM & MFP + Feat & 23.23 $\pm$ 1.08 & 40.07 & 31.12 $\pm$ 1.65 & 52.55 & 0.75 $\pm$ 0.04 & 0.43 & 0.86 $\pm$ 0.02 & 0.66 & 0.85 $\pm$ 0.02 & 0.64 \\
    Catboost & LS + Feat & 26.35 $\pm$ 1.56 & 39.24 & 37.26 $\pm$ 2.51 & 52.45 & 0.62 $\pm$ 0.04 & 0.44 & 0.79 $\pm$ 0.03 & 0.66 & 0.79 $\pm$ 0.02 & \textbf{0.66} \\
    Catboost & MFP + Feat & 21.02 $\pm$ 1.48 & \textbf{39.04} & 31.92 $\pm$ 2.43 & \textbf{52.03} & 0.73 $\pm$ 0.03 & \textbf{0.44} & 0.85 $\pm$ 0.02 & \textbf{0.66} & 0.85 $\pm$ 0.02 & 0.65 \\
    FFNN & LS + Feat & 23.51 $\pm$ 1.43 & 45.15 & 32.09 $\pm$ 2.1 & 58.69 & 0.72 $\pm$ 0.03 & 0.32 & 0.85 $\pm$ 0.02 & 0.57 & 0.84 $\pm$ 0.02 & 0.55 \\
    FFNN & MFP + Feat & 18.19 $\pm$ 1.05 & 42.65 & 25.42 $\pm$ 1.58 & 56.18 & 0.82 $\pm$ 0.02 & 0.36 & 0.91 $\pm$ 0.01 & 0.6 & 0.90 $\pm$ 0.01 & 0.59 \\
    GAT + FFNN & SMILES + Feat & 26.31 $\pm$ 0.37 & 40.7 & 37.13 $\pm$ 0.64 & 54.9 & 0.64 $\pm$ 0.01 & 0.39 & 0.80 $\pm$ 0.00 & 0.63 & 0.80 $\pm$ 0.01 & 0.62 \\
    CNN + FFNN & SMILES + Feat & 26.03 $\pm$ 0.23 & 40.85 & 37.03 $\pm$ 0.49 & 55.6 & 0.65 $\pm$ 0.01 & 0.4 & 0.80 $\pm$ 0.00 & 0.63 & 0.80 $\pm$ 0.00 & 0.63 \\
    LSTM + FFNN & SMILES + Feat & 26.27 $\pm$ 0.42 & 40.7 & 36.87 $\pm$ 0.61 & 55.14 & 0.65 $\pm$ 0.01 & 0.39 & 0.80 $\pm$ 0.00 & 0.62 & 0.799 $\pm$ 0.01 & 0.62 \\
    \hline 
    \end{tabular}
    }
    \caption{Performance comparison of ML models predicting drug response for AML patients. Here we highlight the model with the best prediction metrics on the test set.}
    \label{table:t1}
\end{table}
%%%%%%%%%%%%%%%%%%%%%%%%%%%%%%%%%%%%%%%%
\section{Conclusion}
In this study, we developed PT-AML, a machine learning framework for predicting personalized drug responses in acute myeloid leukemia patients. By leveraging the rich multi-omics and clinical data from the BeatAML cohort, we engineered comprehensive representations of drugs, patient profiles, and biological pathways to accurately estimate drug sensitivity.
Our optimal PT-AML model, based on the CatBoost algorithm, demonstrated strong predictive performance, achieving a mean absolute error of 39.04 and a Pearson correlation of 0.66 on the independent test set. This highlights the model's ability to generalize to unseen patient data.

While our results are promising, further validation on prospective clinical cohorts is necessary to fully assess the model's real-world applicability. Additionally, interpretability analyses using methods like SHAP \cite{lundberg2017unified} could provide valuable insights into the key features driving drug response predictions. In conclusion, PT-AML represents a significant step towards precision medicine in AML treatment. By accurately predicting patient-specific drug responses, this framework has the potential to guide clinicians in selecting optimal therapies, ultimately improving outcomes for AML patients.


\bibliographystyle{unsrt}
\bibliography{mybibliography}

\end{document}
